
\documentclass{article}
\usepackage{colortbl}
\usepackage{makecell}
\usepackage{multirow}
\usepackage{supertabular}

\begin{document}

\newcounter{utterance}

\centering \large Interaction Transcript for game `wordle', experiment `high\_frequency\_words\_no\_clue\_no\_critic', episode 7 with deepseek{-}v4{-}pro.
\vspace{24pt}

{ \footnotesize  \setcounter{utterance}{1}
\setlength{\tabcolsep}{0pt}
\begin{supertabular}{c@{$\;$}|p{.15\linewidth}@{}p{.15\linewidth}p{.15\linewidth}p{.15\linewidth}p{.15\linewidth}p{.15\linewidth}}
    \# & \multicolumn{2}{c}{Player} && \multicolumn{2}{c}{Game Master} \\
    \hline

    \theutterance \stepcounter{utterance}  
    & & & \multicolumn{4}{p{0.6\linewidth}}{
        \cellcolor[rgb]{0.9,0.9,0.9}{
            \makecell[{{p{\linewidth}}}]{
                \texttt{\tiny{[P1$\langle$GM]}}
                \texttt{Ste jezikovni čarovnik, ki rad ugiba besede z uporabo danih pravil.} \\
\\ 
\texttt{Dobrodošli v Wordlu! Imate šest poskusov, da uganete ciljno besedo, veljavno slovensko besedo s petimi malimi črkami. Prosimo, uporabite oznaki "razlaga:" in "ugibanje:", da za vsak poskus podate jedrnato razlago.} \\
\\ 
\texttt{Na primer, če je vaš poskus "otrok", naj bo vaš odgovor} \\
\texttt{razlaga: to je pogosta petčrkovna slovenska beseda in svoj poskus začenjam s to besedo.} \\
\texttt{ugibanje: otrok} \\
\\ 
\texttt{Po vsakem poskusu bo vaš odgovor preverjen in prejeli boste povratne informacije, ki bodo pokazale, katere črke so pravilne (zelena), katere črke so pravilne, vendar na napačnem mestu (rumena), in katere črke so napačne (rdeča). Te povratne informacije so lahko koristne pri določanju, katere črke vključiti ali izključiti v naslednjem poskusu.} \\
\\ 
\texttt{Na primer, povratne informacije za "otrok" bi lahko bile:} \\
\texttt{povratna\_informacija\_ugibanja: o<rumena> t<rumena> r<zelena> o<rumena> k<rdeča>} \\
\\ 
\texttt{Razlaga naj vsebuje podrobnosti o tem, kako se povratna\_informacija\_ugibanja uporabi za oblikovanje novega poskusa.} \\
\\ 
\texttt{Ne ustvarjajte nobenega drugega besedila.} \\
\\ 
\texttt{Začnimo z vašim prvim poskusom.} \\
            }
        }
    }
     \\ \\

    \theutterance \stepcounter{utterance}  
    & \multicolumn{4}{p{0.6\linewidth}}{
        \cellcolor[rgb]{0.561,0.737,0.561}{
            \makecell[{{p{\linewidth}}}]{
                \texttt{\tiny{[P1$\rangle$GM]}}
                \texttt{razlaga: Za prvi poskus bom uporabil besedo "kamen", ker vsebuje pogoste soglasnike in samoglasnik.} \\
\texttt{ugibanje: kamen} \\
            }
        }
    }
    & & \\ \\

    \theutterance \stepcounter{utterance}  
    & & & \multicolumn{2}{p{0.3\linewidth}}{
        \cellcolor[rgb]{0.9,0.9,0.9}{
            \makecell[{{p{\linewidth}}}]{
                \texttt{\tiny{[GM$|$GM]}}
                \texttt{attempts = 1} \\
\texttt{target\_word = glava} \\
\texttt{guess = kamen} \\
\texttt{guess\_feedback = k<red> a<yellow> m<red> e<red> n<red>} \\
            }
        }
    }
    & & \\ \\

    \theutterance \stepcounter{utterance}  
    & & & \multicolumn{4}{p{0.6\linewidth}}{
        \cellcolor[rgb]{0.9,0.9,0.9}{
            \makecell[{{p{\linewidth}}}]{
                \texttt{\tiny{[P1$\langle$GM]}}
                \texttt{povratna\_informacija\_ugibanja: k<rdeča> a<rumena> m<rdeča> e<rdeča> n<rdeča>} \\
\\ 
\texttt{Odgovorite samo v tej obliki.} \\
\texttt{razlaga: podrobnosti} \\
\texttt{ugibanje: beseda} \\
            }
        }
    }
     \\ \\

    \theutterance \stepcounter{utterance}  
    & \multicolumn{4}{p{0.6\linewidth}}{
        \cellcolor[rgb]{0.561,0.737,0.561}{
            \makecell[{{p{\linewidth}}}]{
                \texttt{\tiny{[P1$\rangle$GM]}}
                \texttt{razlaga: Črka 'a' je v besedi, vendar ne na drugem mestu. 'K', 'm', 'e' in 'n' so izključene. Poskusil bom besedo "slika", ki vsebuje 'a' na drugem mestu in preverja nove pogoste črke.} \\
\texttt{ugibanje: slika} \\
            }
        }
    }
    & & \\ \\

    \theutterance \stepcounter{utterance}  
    & & & \multicolumn{2}{p{0.3\linewidth}}{
        \cellcolor[rgb]{0.9,0.9,0.9}{
            \makecell[{{p{\linewidth}}}]{
                \texttt{\tiny{[GM$|$GM]}}
                \texttt{attempts = 2} \\
\texttt{target\_word = glava} \\
\texttt{guess = slika} \\
\texttt{guess\_feedback = s<red> l<green> i<red> k<red> a<green>} \\
            }
        }
    }
    & & \\ \\

    \theutterance \stepcounter{utterance}  
    & & & \multicolumn{4}{p{0.6\linewidth}}{
        \cellcolor[rgb]{0.9,0.9,0.9}{
            \makecell[{{p{\linewidth}}}]{
                \texttt{\tiny{[P1$\langle$GM]}}
                \texttt{povratna\_informacija\_ugibanja: s<rdeča> l<zelena> i<rdeča> k<rdeča> a<zelena>} \\
\\ 
\texttt{Odgovorite samo v tej obliki.} \\
\texttt{razlaga: podrobnosti} \\
\texttt{ugibanje: beseda} \\
            }
        }
    }
     \\ \\

    \theutterance \stepcounter{utterance}  
    & \multicolumn{4}{p{0.6\linewidth}}{
        \cellcolor[rgb]{0.561,0.737,0.561}{
            \makecell[{{p{\linewidth}}}]{
                \texttt{\tiny{[P1$\rangle$GM]}}
                \texttt{razlaga: Črki 'l' in 'a' sta na pravilnih mestih (drugo in peto). 'S', 'i', 'k' so izključene. Beseda se torej začne z neznano črko, sledi 'l', nato dve neznani črki in konča z 'a'. Poskusil bom besedo "vlaga", ki ustreza vzorcu \_ l \_ \_ a in preverja nove črke.} \\
\texttt{ugibanje: vlaga} \\
            }
        }
    }
    & & \\ \\

    \theutterance \stepcounter{utterance}  
    & & & \multicolumn{2}{p{0.3\linewidth}}{
        \cellcolor[rgb]{0.9,0.9,0.9}{
            \makecell[{{p{\linewidth}}}]{
                \texttt{\tiny{[GM$|$GM]}}
                \texttt{attempts = 3} \\
\texttt{target\_word = glava} \\
\texttt{guess = vlaga} \\
\texttt{guess\_feedback = v<yellow> l<green> a<green> g<yellow> a<green>} \\
            }
        }
    }
    & & \\ \\

    \theutterance \stepcounter{utterance}  
    & & & \multicolumn{4}{p{0.6\linewidth}}{
        \cellcolor[rgb]{0.9,0.9,0.9}{
            \makecell[{{p{\linewidth}}}]{
                \texttt{\tiny{[P1$\langle$GM]}}
                \texttt{povratna\_informacija\_ugibanja: v<rumena> l<zelena> a<zelena> g<rumena> a<zelena>} \\
\\ 
\texttt{Odgovorite samo v tej obliki.} \\
\texttt{razlaga: podrobnosti} \\
\texttt{ugibanje: beseda} \\
            }
        }
    }
     \\ \\

    \theutterance \stepcounter{utterance}  
    & \multicolumn{4}{p{0.6\linewidth}}{
        \cellcolor[rgb]{0.561,0.737,0.561}{
            \makecell[{{p{\linewidth}}}]{
                \texttt{\tiny{[P1$\rangle$GM]}}
                \texttt{razlaga: 'l', 'a' (tretje mesto) in 'a' (peto mesto) so zelene. 'V' in 'g' sta v besedi, vendar na napačnih mestih. Prvo mesto je torej 'v' ali 'g' (ali kakšna druga črka), četrto mesto pa preostala od teh dveh. Beseda ima obliko {[}\_{]} l a {[}g/v{]} a. Poskusil bom "blaga", ker ustreza vzorcu in preveri novo črko 'b' na prvem mestu ter postavi 'g' na četrto.} \\
\texttt{ugibanje: blaga} \\
            }
        }
    }
    & & \\ \\

    \theutterance \stepcounter{utterance}  
    & & & \multicolumn{2}{p{0.3\linewidth}}{
        \cellcolor[rgb]{0.9,0.9,0.9}{
            \makecell[{{p{\linewidth}}}]{
                \texttt{\tiny{[GM$|$GM]}}
                \texttt{attempts = 4} \\
\texttt{target\_word = glava} \\
\texttt{guess = blaga} \\
\texttt{guess\_feedback = b<red> l<green> a<green> g<yellow> a<green>} \\
            }
        }
    }
    & & \\ \\

    \theutterance \stepcounter{utterance}  
    & & & \multicolumn{4}{p{0.6\linewidth}}{
        \cellcolor[rgb]{0.9,0.9,0.9}{
            \makecell[{{p{\linewidth}}}]{
                \texttt{\tiny{[P1$\langle$GM]}}
                \texttt{povratna\_informacija\_ugibanja: b<rdeča> l<zelena> a<zelena> g<rumena> a<zelena>} \\
\\ 
\texttt{Odgovorite samo v tej obliki.} \\
\texttt{razlaga: podrobnosti} \\
\texttt{ugibanje: beseda} \\
            }
        }
    }
     \\ \\

    \theutterance \stepcounter{utterance}  
    & \multicolumn{4}{p{0.6\linewidth}}{
        \cellcolor[rgb]{0.561,0.737,0.561}{
            \makecell[{{p{\linewidth}}}]{
                \texttt{\tiny{[P1$\rangle$GM]}}
                \texttt{razlaga: 'b' je izključena. 'g' je še vedno rumena na četrtem mestu, torej mora biti 'g' na prvem mestu, ker je to edino preostalo prosto mesto. Beseda se torej glasi g l a \_ a. Na četrtem mestu manjka črka, ki mora biti prejšnja rumena črka 'v' (iz prejšnjega poskusa), saj druge možnosti ni. Poskusil bom "glava".} \\
\texttt{ugibanje: glava} \\
            }
        }
    }
    & & \\ \\

    \theutterance \stepcounter{utterance}  
    & & & \multicolumn{2}{p{0.3\linewidth}}{
        \cellcolor[rgb]{0.9,0.9,0.9}{
            \makecell[{{p{\linewidth}}}]{
                \texttt{\tiny{[GM$|$GM]}}
                \texttt{attempts = 5} \\
\texttt{target\_word = glava} \\
\texttt{guess = glava} \\
\texttt{guess\_feedback = g<green> l<green> a<green> v<green> a<green>} \\
            }
        }
    }
    & & \\ \\

    \theutterance \stepcounter{utterance}  
    & & & \multicolumn{2}{p{0.3\linewidth}}{
        \cellcolor[rgb]{0.9,0.9,0.9}{
            \makecell[{{p{\linewidth}}}]{
                \texttt{\tiny{[GM$|$GM]}}
                \texttt{game\_result = WIN} \\
            }
        }
    }
    & & \\ \\

\end{supertabular}
}

\end{document}
